\customchapter{ABSTRACT}

\begin{center}
\large \vspace{-1.5cm} \textbf{CAVA: TOPIC SEGMENTATION OF LONG VIDEO BY CONSENSUS-BASED MULTIMODAL AGGREGATION WITHOUT REPRESENTATION FUSION}
\end{center}

Automatic analysis of long video requires splitting it into meaningful units, and no
modality should hide the others while doing so. This thesis builds on a diagnosis from
previous work: in A2Summ-X, an extension of the A2Summ summarization model
\cite{he2023a2summ}, fusing modalities into a single representation lets text dominate the
weighting between modalities (weights of 0.75 to 0.89) and cancels the visual contribution
before any decision is made. We propose CAVA (Consumer-Adaptive Video Architecture), a topic
segmentation system for long video that does not fuse representations. Each signal
(transcript, image, voice and the descriptions of a vision--language model) emits its own
boundaries and topic labels. Topics emerge from a quorum vote, signals without evidence in a
stretch abstain, and every group records which signals support it.

The system produces a topic map at two granularity levels that links non-adjacent segments
when a topic returns. Five version formats are derived from the map: chapters, a 60-second
recap, 10\,\% highlights, topic threads and speaker threads. Validation uses pre-registered
gates on unseen data. On the 47 films of the MovieNet test split \cite{huang2020movienet},
the map with abstention beats concatenation fusion ($z = 2.46$) and the contiguous manager
($z = 2.48$) in pairwise MCC. The gain depends on abstention: without it, consensus is
indistinguishable from fusion ($z = 0.80$), and the visual signal alone scores a higher MCC
than any combination. Versions meet their budget in every content type evaluated (film,
meeting, web and user video), and speaker identity reaches an ARI of 0.73 on AMI
\cite{carletta2007ami}. Three gates are not met: highlight quality against simple baselines
(EXP-031/037/038), non-adjacent links judged as the same plot (G-enlace-2), and narrator
detection in documentaries (EXP-039). Of the 26 modules of the frozen system, 10 are
validated. The human evaluation (EXP-040) is designed, with its set of 9 videos already
processed, and is pending pre-registration and execution.

\vspace{0.5em}
\noindent \textbf{Keywords:}\\
\noindent Video topic segmentation; multimodal aggregation; consensus voting;
vision--language models; consumer-adaptive versions.
